\documentclass[10pt,a4paper]{article}

\usepackage[margin=1.6cm,top=1.35cm,bottom=1.45cm]{geometry}
\usepackage[T1]{fontenc}
\usepackage[utf8]{inputenc}
\usepackage[default]{sourcesanspro}
\usepackage[scaled=0.88]{inconsolata}
\usepackage{microtype}
\usepackage{enumitem}
\usepackage{tabularx}
\usepackage{array}
\usepackage{xcolor}
\usepackage{hyperref}
\usepackage{xurl}
\usepackage{fontawesome5}
\usepackage{fancyhdr}
\usepackage{lastpage}
\usepackage{needspace}
\usepackage{ragged2e}
\usepackage[most]{tcolorbox}

\definecolor{Ink}{HTML}{132534}
\definecolor{Navy}{HTML}{173D55}
\definecolor{Teal}{HTML}{087E8B}
\definecolor{Muted}{HTML}{5D6B75}
\definecolor{Line}{HTML}{D8E2E7}
\definecolor{Pale}{HTML}{EEF7F8}
\definecolor{Warm}{HTML}{F5F1E8}

\hypersetup{
  colorlinks=true,
  urlcolor=Teal,
  linkcolor=Teal,
  pdfauthor={Tianjian Qin},
  pdftitle={Curriculum Vitae - Tianjian Qin},
  pdfsubject={Data science, computational biology, and health research}
}

\setlength{\parindent}{0pt}
\setlength{\parskip}{0pt}
\setlength{\tabcolsep}{0pt}
\setlength{\footskip}{0.72cm}
\renewcommand{\familydefault}{\sfdefault}
\emergencystretch=2em
\hyphenpenalty=10000
\exhyphenpenalty=10000

\pagestyle{fancy}
\fancyhf{}
\renewcommand{\headrulewidth}{0pt}
\fancyfoot[L]{\scriptsize\color{Muted} Tianjian Qin \enspace\textcolor{Line}{|}\enspace Curriculum Vitae}
\fancyfoot[R]{\scriptsize\color{Muted} \thepage\ /\ \pageref*{LastPage}}

\newlist{cvlist}{itemize}{1}
\setlist[cvlist]{
  label=\raisebox{0.2ex}{\scriptsize\color{Teal}\faIcon{angle-right}},
  leftmargin=1.15em,
  rightmargin=0pt,
  topsep=0.25em,
  itemsep=0.12em,
  parsep=0pt,
  partopsep=0pt
}

\newlist{publist}{enumerate}{1}
\setlist[publist]{
  label=\textcolor{Teal}{\scriptsize\bfseries\arabic*.},
  leftmargin=1.55em,
  labelsep=0.45em,
  topsep=0.2em,
  itemsep=0.36em,
  parsep=0pt,
  partopsep=0pt
}

\newcommand{\cvsection}[1]{%
  \Needspace{4\baselineskip}%
  \vspace{0.72em}%
  \noindent
  \begin{tabularx}{\textwidth}{@{}l@{\hspace{0.58em}}X@{}}
    {\color{Teal}\ttfamily\bfseries //} &
    {\large\bfseries\color{Teal} #1}
  \end{tabularx}
  \vspace{0.16em}\par
  {\color{Line}\hrule height 0.7pt}
  \vspace{0.42em}
}

\newcommand{\cvsubhead}[1]{%
  \Needspace{2.5\baselineskip}%
  \vspace{0.45em}%
  {\small\bfseries\color{Teal}\MakeUppercase{#1}}\par
  \vspace{0.18em}
}

\newcommand{\cventry}[4]{%
  \Needspace{9\baselineskip}%
  \begin{tabularx}{\textwidth}{@{}X>{\raggedleft\arraybackslash}p{3.05cm}@{}}
    {\bfseries\color{Ink} #1} & {\bfseries\color{Teal} #2} \\
    {\small\color{Muted} #3} & {}
  \end{tabularx}
  \vspace{0.12em}
  #4
  \vspace{0.48em}
}

\newcommand{\statusbadge}[2]{%
  \begingroup
  \setlength{\fboxsep}{3.2pt}%
  \colorbox{#1}{\textcolor{Teal}{\scriptsize\bfseries\strut #2}}%
  \endgroup
}

\newcommand{\fundingentry}[5]{%
  \Needspace{5\baselineskip}%
  \begin{tabularx}{\textwidth}{@{}X>{\raggedleft\arraybackslash}p{3.2cm}@{}}
    {\bfseries\color{Ink} #1} & #2 \\
    {\small\color{Muted} #3} & {\small\color{Muted} #4}
  \end{tabularx}
  \vspace{0.12em}
  #5
  \vspace{0.48em}
}

\newcommand{\projectentry}[4]{%
  \Needspace{5\baselineskip}%
  \begin{tabularx}{\textwidth}{@{}X@{\hspace{0.3cm}}>{\raggedleft\arraybackslash}p{3.05cm}@{}}
    {\bfseries\color{Ink} #1} & {\small\bfseries\color{Teal} #2} \\
    {\small\color{Muted} #3} & {}
  \end{tabularx}
  \vspace{0.12em}
  #4
  \vspace{0.48em}
}

\newtcolorbox{featurebox}{
  colback=Pale,
  colframe=Pale,
  boxrule=0pt,
  sharp corners,
  left=2.4mm,
  right=2.4mm,
  top=1.8mm,
  bottom=1.8mm,
  before skip=0.15em,
  after skip=0.55em
}

\newcommand{\theme}[2]{%
  {\scriptsize\bfseries\color{Teal}\MakeUppercase{#1}}\par
  \vspace{0.05em}
  {\small\color{Ink}#2}
}

\newcommand{\tool}[1]{\texttt{\color{Navy}#1}}
\newcommand{\sep}{\enspace\textcolor{Line}{|}\enspace}

\begin{document}
\RaggedRight

\begin{minipage}[t]{0.54\textwidth}
  \vspace{0pt}
  {\fontsize{28}{30}\selectfont\bfseries\color{Navy} Tianjian Qin, PhD}\par
  \vspace{0.24em}
  {\large\color{Teal} Computational Biologist \& Data Scientist}\par
\end{minipage}
\hfill
\begin{minipage}[t]{0.42\textwidth}
  \vspace{0.15em}
  \raggedleft\small\color{Muted}
  \faIcon{envelope}\enspace
  \href{mailto:tianjian.qin@wur.nl}{tianjian.qin@wur.nl}
  \sep
  \faIcon{map-marker-alt}\enspace Netherlands\par
  \vspace{0.22em}
  \faIcon{globe}\enspace\href{https://qtj.me}{qtj.me}
  \sep
  \faIcon{orcid}\enspace\href{https://orcid.org/0000-0003-2540-1778}{0000-0003-2540-1778}\par
  \vspace{0.22em}
  \faIcon{project-diagram}\enspace\href{https://herdlink.nl}{HerdLink.nl}
  \sep
  \faIcon{github}\enspace\href{https://github.com/EvoLandEco}{EvoLandEco}
  \sep
  \faIcon{linkedin}\enspace\href{https://www.linkedin.com/in/tjqin}{tjqin}
\end{minipage}

\vspace{0.75em}
{\color{Teal}\hrule height 1.3pt}

\cvsection{Research Profile}

\begin{cvlist}
    \item I develop machine learning and deep learning methods, statistical models, simulation
    frameworks, and research software for complex biological systems. My work spans evolutionary
    biology, ecology, infectious disease epidemiology, and One Health research. I work with domain
    experts to turn scientific questions into transparent models, scalable pipelines, and usable tools.
\end{cvlist}

\cvsection{Research Experience}

\cventry
  {Postdoctoral Researcher}
  {2025--present}
  {Wageningen University \& Research, Infectious Disease Epidemiology Group}
  {%
    \begin{cvlist}
      \item Develop generative and predictive models for dynamic livestock movement networks,
      including neural models for spatial point patterns, Transformer models for link forecasting,
      stochastic block models, and transmission simulations.
      \item Build reproducible HPC pipelines that integrate movement records, farm metadata,
      spatial information, temporal covariates, and epidemiological context.
      \item Translate empirical data and stakeholder questions into modelling workflows, open software,
      visual analysis, and decision support outputs across animal and public health domains.
    \end{cvlist}
  }

\cventry
  {Doctoral Researcher in Theoretical \& Computational Biology}
  {2019--2026}
  {University of Groningen, Groningen Institute for Evolutionary Life Sciences}
  {%
    \begin{cvlist}
      \item Developed the EVE stochastic birth-death model and high-performance simulation and
      inference methods for diversification dependent on evolutionary relatedness.
      \item Designed EvoNN, an ensemble learning framework combining graph, recurrent, and dense
      neural networks for parameter inference from phylogenetic trees.
      \item Built and maintained reproducible research software in \tool{R}, \tool{C++}, and
      \tool{Python}, with extensive CPU and GPU computing, testing, documentation, and open releases.
    \end{cvlist}
  }

\cventry
  {MSc Researcher in Wetland Ecology \& Invasion Biology}
  {2016--2019}
  {Beijing Forestry University, School of Nature Conservation}
  {%
    \begin{cvlist}
      \item Conducted field surveys across 18 Chinese provinces and received training in controlled
      experiments, GIS, phylogenetic metrics, community and trait analyses, spatial analysis, and
      environmental data analysis.
      \item Built automated image analysis and phylogenetic reconstruction workflows and published
      quantitative studies of invasion resistance and community assembly.
    \end{cvlist}
  }

\cvsection{Research Funding \& Proposal Development}

\fundingentry
  {NextdAI: Integrated AI-Driven Early Detection System}
  {\statusbadge{Warm}{IN DEVELOPMENT}}
  {NWO KIC Early Detection of (Emerging) Infectious Disease Outbreaks}
  {2026}
  {\begin{cvlist}
    \item Co-developing a novel AI architecture and a data infrastructure plan for heterogeneous
      surveillance streams, anomaly detection, neuro-symbolic characterisation,
      forecasting, and knowledge-based decision support.
  \end{cvlist}}

\fundingentry
  {SSS-mod: Preparedness Package for Future Respiratory Outbreaks}
  {\statusbadge{Pale}{AWARDED}}
  {ZonMw Implementatie-impuls COVID-19 programma}
  {Mar--Dec 2026}
  {\begin{cvlist}
    \item Project team member. Responsible for programming for data preparation,
      recalibration, model output generation, and interface development, supporting the reuse of a
      setting-specific transmission model during future respiratory outbreaks.
  \end{cvlist}}

\fundingentry
  {Small Compute Applications}
  {\statusbadge{Pale}{AWARDED}}
  {NWO \& SURF}
  {2025}
  {\begin{cvlist}
    \item Compute allocation supporting data-intensive simulation, neural inference, and model
      evaluation on high-performance computing systems.
  \end{cvlist}}

\cvsection{Selected Research Projects}

\begin{featurebox}
  \begin{tabularx}{\linewidth}{@{}X>{\raggedleft\arraybackslash}p{3.4cm}@{}}
    {\large\bfseries\color{Navy} HerdLink.nl} &
    {\statusbadge{white}{LIVE PLATFORM}} \\
    {\small\color{Muted} Developer \sep \tool{JavaScript} \sep \tool{D3.js}} &
    {\small\href{https://herdlink.nl}{https://herdlink.nl}}
  \end{tabularx}
  \vspace{0.25em}

  Developed and deployed a public web platform for exploring anonymised and aggregated Dutch
  livestock trade networks. HerdLink links graph and map views with node focus, temporal replay,
  network summaries, and transmission simulation for research exploration, communication, and education.
\end{featurebox}

\projectentry
  {EUPAHW HEV: transmission pathways across pigs, wild boar, and humans}
  {Researcher}
  {European Partnership on Animal Health and Welfare}
  {\begin{cvlist}
    \item Developing research questions and analysis strategies that join phylodynamics,
      animal movement networks, wildlife and landscape information, and intervention timing. The work
      targets transmission direction, spillover between species, long-range lineage movement, and
      generalisable One Health surveillance design.
  \end{cvlist}}

\projectentry
  {NetForge: dynamic livestock network modelling}
  {Developer}
  {\tool{Python} \sep network science \sep epidemic simulation}
  {\begin{cvlist}
    \item Developing a research toolkit for movement network synthesis, statistical evaluation, metadata
      generation, and transmission experiments. Current work includes physics-informed network
      generators and models for spatial and temporal link prediction.
  \end{cvlist}}

\projectentry
  {EVE \& EvoNN: simulation and neural inference from phylogenies}
  {Researcher}
  {\tool{R} \sep \tool{C++} \sep \tool{Python} \sep GPU computing}
  {\begin{cvlist}
    \item Created stochastic models, simulators, and ensemble neural networks that learn jointly from graph
      structure, branching times, and summary statistics. The work supports peer-reviewed publications,
      reusable packages, and a completed doctoral thesis.
  \end{cvlist}}

\cvsection{Teaching \& Mentoring}

\begin{cvlist}
  \item Teaching assistant for annual coding assignments and computer practicals in island
  biogeography at the University of Groningen.
  \item Guest lectures, tutorials, and seminars on stochastic processes, diversification models,
  deep learning, neural networks, and network statistics.
  \item Lectures and workshops at Wageningen University \& Research and Wageningen Bioveterinary
  Research on phylodynamics, birth-death models, and network modelling.
  \Needspace{4\baselineskip}
  \item Pair programming sessions on test-driven development and collaborative software design.
  \item Informal mentoring of BSc and MSc students in biostatistics, data visualisation, and
  scientific writing, and of junior PhD researchers in planning and day-to-day research practice.
\end{cvlist}

\Needspace{15\baselineskip}

\cvsection{Education}

\cventry
  {PhD in Theoretical Biology}
  {2019--2026}
  {University of Groningen, Netherlands}
  {\begin{cvlist}
      \item Thesis: \emph{Diversification Models and Neural Inference}. \sep
      \href{https://research.rug.nl/files/1515277445/Complete_thesis.pdf}
      {\faIcon{download}\enspace Thesis PDF}
      \item Supervisors: Prof. Dr.\ Rampal Etienne, Prof. Dr.\ Luis Valente, and Dr.\ Koen van Benthem.
  \end{cvlist}}

\cventry
  {MSc in Wetland Ecology}
  {2016--2019}
  {Beijing Forestry University, China}
  {\begin{cvlist}
      \item Supervisors: Prof. Dr.\ Hongli Li and Prof. Dr.\ Feihai Yu.
  \end{cvlist}}

\cventry
  {BSc in Marine Biology}
  {2012--2016}
  {Nanjing Normal University, China}
  {}

\cvsection{Publications}

\cvsubhead{Peer-reviewed articles}

\begin{publist}
  \item \textbf{Qin, T.}, van Benthem, K., Valente, L.\textsuperscript{\dag},
  \& Etienne, R. S.\textsuperscript{\dag} (2026).
  Parameter estimation from phylogenetic trees using neural networks and ensemble learning.
  \emph{Systematic Biology}, 75(2), 344--365.

  \item \textbf{Qin, T.}, Valente, L.\textsuperscript{\dag},
  \& Etienne, R. S.\textsuperscript{\dag} (2025).
  Impact of evolutionary relatedness on species diversification and tree shape.
  \emph{Journal of Theoretical Biology}, 598, 111992.

  \item Sun, K., Liu, X.-S., \textbf{Qin, T.-J.}, Jiang, F., Cai, J.-F.,
  Shen, Y.-L., A, S.-H., \& Li, H.-L. (2021).
  Relative abundance of invasive plants more effectively explains the response of wetland
  communities to different invasion degrees than phylogenetic evenness.
  \emph{Journal of Plant Ecology}.

  \item \textbf{Qin, T.}, Zhou, J., Sun, Y., M\"uller-Sch\"arer, H., Luo, F.,
  Dong, B., Li, H., \& Yu, F.-H. (2020).
  Phylogenetic diversity is a better predictor of wetland community resistance to
  \emph{Alternanthera philoxeroides} invasion than species richness.
  \emph{Plant Biology}.

  \item Wan, J.-Z., Wang, M.-Z., \textbf{Qin, T.}, Bu, X.-Q., Li, H.-L.,
  \& Yu, F.-H. (2019).
  Spatial environmental heterogeneity may be the driver of functional trait variation in
  \emph{Hydrocotyle vulgaris}, an aquatic plant invader.
  \emph{Aquatic Biology}.

  \item \textbf{Qin, T.-J.}*, Guan, Y.-T.*, Quan, H., Dong, B.-C., Luo, F.-L.,
  Zhang, M.-X., Li, H.-L., \& Yu, F.-H. (2019).
  Growth traits of the exotic plant \emph{Hydrocotyle vulgaris} and the evenness of resident
  plant communities are mediated by community age, not species diversity.
  \emph{Weed Research}.

  \item \textbf{Qin, T.-J.}, Guan, Y.-T., Zhang, M.-X., Li, H.-L.,
  \& Yu, F.-H. (2018).
  Sediment type and nitrogen deposition affect the relationship between
  \emph{Alternanthera philoxeroides} and experimental wetland plant communities.
  \emph{Marine and Freshwater Research}.

  \item Liu, L., Guan, Y.-T., \textbf{Qin, T.-J.}, Wang, Y.-Y., Li, H.-L.,
  \& Zhi, Y.-B. (2018).
  Effects of water regime on the growth of the submerged macrophyte
  \emph{Ceratophyllum demersum} at different densities.
  \emph{Journal of Freshwater Ecology}.
\end{publist}

\cvsubhead{Preprint}

\begin{publist}
    \item \textbf{Qin, T.}, van Benthem, K., Valente, L.\textsuperscript{\dag},
    \& Etienne, R. S.\textsuperscript{\dag} (2026).
    Identifying evolutionary relatedness effects on diversification from phylogenies using neural
    networks. \emph{bioRxiv}.
\end{publist}

\cvsubhead{Conference}

\begin{publist}
    \item \textbf{Qin, T.}, Atamer Balkan, B., Schmid, B., \& ten Bosch, Q.
    Generating livestock movement networks that preserve epidemic behaviour with physics informed
    machine learning. \emph{ModAH Conference 2026, Nantes, France}.
\end{publist}

\vspace{0.35em}
{\small\color{Muted}* Equal contribution. \textsuperscript{\dag} Joint senior authors.}

\cvsection{Honours \& Awards}

\begin{cvlist}
  \item \textbf{\href{https://www.sciencedirect.com/journal/journal-of-theoretical-biology/about/news/the-winner-of-the-denise-kirschner-best-student-paper-prize}
  {Denise Kirschner Best Student Paper Prize \faIcon{external-link-alt}}}, \emph{Journal of Theoretical Biology}
  (2025). Inaugural prize for graduate student research.
  \item \textbf{First-Class Academic Performance Award}, Beijing Forestry University
  (2017--2018 and 2018--2019).
  \item \textbf{Graduate Academic Innovation Award}, Beijing Forestry University (2018).
\end{cvlist}

\cvsection{Academic \& Community Service}

\begin{cvlist}
  \item \textbf{Peer reviewer} for \emph{Systematic Biology}, \emph{PLOS Computational Biology},
  \emph{Epidemics}, \emph{BMC Public Health}, and \emph{BMC Plant Biology}.
  \item \textbf{Contributor} to various GitHub open-source repositories.
  \item \textbf{Vice President}, Graduate Student Council, School of Nature Conservation, Beijing Forestry
  University (2017--2018).
  \item \textbf{Student representative}, Wetland Ecology Group, Beijing Forestry University (2016--2019).
\end{cvlist}

\Needspace{8\baselineskip}
\cvsection{Research Software}

\projectentry
  {\href{https://qtj.me/NetForge}{NetForge}}
  {Developer}
  {Movement network modelling and transmission simulation \sep \tool{Python}}
  {\begin{cvlist}
      \item Command-line tools for network synthesis, metadata generation, evaluation, and epidemic
      experiments.
  \end{cvlist}}

\projectentry
  {\href{https://github.com/EvoLandEco/EvoNN}{EvoNN} \&
  \href{https://github.com/EvoLandEco/evesim}{evesim}}
  {Developer}
  {Neural inference and stochastic simulation \sep \tool{R}/\tool{Python}/\tool{C++}}
  {\begin{cvlist}
      \item Pretrained neural inference from phylogenies and fast simulation of evolutionary relatedness
      dependent birth-death processes.
  \end{cvlist}}

\projectentry
  {\href{https://github.com/EvoLandEco/NetSpectra}{NetSpectra},
  \href{https://github.com/EvoLandEco/EvoLab}{EvoLab}, \&
  \href{https://github.com/EvoLandEco/miniape}{miniape}}
  {Developer}
  {Interactive scientific tools \sep \tool{JavaScript}}
  {\begin{cvlist}
      \item Network dynamics, lineage simulation, tree reconstruction, and browser-based phylogenetic data
      handling.
  \end{cvlist}}

\projectentry
  {\href{https://github.com/thijsjanzen/treestats}{treestats},
  \href{https://github.com/rsetienne/DDD}{DDD}, \&
  \href{https://github.com/rsetienne/DAISIE}{DAISIE}}
  {Collaborator}
  {High-performance research packages \sep \tool{R}/\tool{C++}}
  {\begin{cvlist}
      \item Contributions to data structures, phylogenetic summary statistics, and simulation and
      inference for diversification and island biogeography models.
  \end{cvlist}}

\cvsection{Technical Expertise}

\begin{tabularx}{\textwidth}{@{}>{\raggedright\arraybackslash}p{3.45cm}@{\hspace{0.42cm}}X@{}}
  \textbf{\color{Navy}Machine learning} &
  Neural networks, transformer models, ensemble learning, and benchmarking. \\[0.34em]

  \textbf{\color{Navy}Statistics \& inference} &
  Stochastic processes, simulation-based inference, and likelihood methods. \\[0.34em]

  \textbf{\color{Navy}Programming} &
  \tool{Python}, \tool{R}, \tool{C/C++}, \tool{JavaScript}, \tool{SQL}; package design, analysis pipelines, and interactive visualisation. \\[0.34em]

  \textbf{\color{Navy}Research computing} &
  Linux, Bash, HPC, data infrastructure, Docker, Git, GitHub workflows, agentic coding, and CI/CD. \\[0.34em]

  \textbf{\color{Navy}Communication} &
  Scientific writing, \LaTeX, TikZ, PGFPlots, D3.js, and interface design for domain experts and students.
\end{tabularx}

\cvsection{Languages}

{\small
\textbf{\color{Navy}Chinese} Native \sep
\textbf{\color{Navy}English} Proficient \sep
\textbf{\color{Navy}Dutch} Basic \sep
\textbf{\color{Navy}German} Basic}

\end{document}
